% !TEX root = Main.tex
%
%\subsection{Lower Bound for the Stochastic \& Adversarial Rewards}
%\label{s:LowBOB}
\paragraph{Why is the best of both worlds unachievable?}
% 
We define a new notion of complexity, $\complexityBoth$ as
%
\[
%
\complexityBoth 
\triangleq
\frac{1}{
\gap_{(1)}}
\max_{k\in\setArms}\frac{k}{\gap_{(k)}}\cdot
%
\]
%
The complexity $\complexityBoth$ captures the difficulty of
stochastic best-arm identification when the learner
must also remain robust to adversarial rewards.
As we detail  in Remark~\ref{rem:LOWbobCOMP}, 
$	\complexitySR\leq\complexityBoth\leq\complexityUnif.$
The gap between $\complexitySR$ and $\complexityBoth$ reflects the
additional difficulty imposed by adversarial robustness.
%Let $\mathcal S$ be the set of stochastic problem such that
%The upcoming Theorem~\ref{th:lowBOB} shows that
%for any $\nArms$ specified gaps, $0<\gap_1=\gap_2\leq\gap_3\leq \ldots\leq \gap_\nArms\leq 1/8$, 
%a stochastic problem can be built with a very similar gap landscape 
%and a complexity changed by at most a factor of $2$ such that on this problem
% it is not possible to have a probability of error 
%smaller than 
%$	\frac{1}{64}\exp\left(-512\frac{\timeHorizon}{\complexityBoth}\right)$
% %on any stochastic problem
%  without  having 
%  a probability of error  not converging to $0$ with large $\timeHorizon$
%   on some adversarial problem with
%  again a similar gap landscape.
%
\begin{restatable}[\textcolor{titleTh}{Lower bound for the BOB challenge}]{theorem}{thi}\label{th:lowBOB}
Fix $\nArms\geq 2$ and reference gaps
$0<\gap_1=\gap_2\leq\cdots\leq\gap_\nArms\leq 1/8$,
and let $\Gaps_4$ be the corresponding problem class.
For any learner and any positive integer horizon $\timeHorizon$
satisfying
\[
\nArms\exp\left(-\frac{\gap_1^2\timeHorizon}{128}\right)
\leq \frac{1}{32},
\]
there exists an i.i.d.\ stochastic problem $\stoPro\in\Gaps_4$
such that, if
\[
\ProErr_{\stoPro}(\timeHorizon)
\leq \frac{1}{64}\exp\left(
-\frac{2048\timeHorizon}{\complexityBoth^{\stoPro}}
\right),
\]
then there exists a deterministic gain table
$\gainVector\in[0,1]^{\nArms\times\timeHorizon}$
belonging to $\Gaps_4$ for which
\[
\ProErr_{\advPro(\gainVector)}(\timeHorizon)
\geq \frac{1}{16}.
\]
Here $\complexityBoth^{\stoPro}$ is the complexity computed
from the gaps of the constructed stochastic problem.
The constructed stochastic problem and gain table may depend
on the learner and on $\timeHorizon$.
      \end{restatable}    
%
\begin{remark}\label{rem:LOWbobCOMP}
	In general, we have
	%
	\[
	\complexitySR\leq\complexityBoth\leq\complexityUnif.
	\]
	%
%We first compare the complexities in a general regime. Then we spell out three different regimes to intuitively explore when the inequalities in the previous equation are strict or not.
Below, we compare the three  complexities in 
three specific gap regimes in order to intuitively explore   whether
the inequalities in the previous equation are strict or not. 
Interestingly, while in two regimes  $\complexitySR=\complexityBoth$, in the third regime, called the `square-root gaps', we can obtain $\complexityBoth=\sqrt{\nArms/2}\complexitySR$.  
This equality shows that on some problems and for large values of $\nArms$,
% 	\todovout{actually really large value on K , like millions !!! :(, i am not sure we want to comment on that no?those numerical constants!}
 our lower bound on 
the complexity of the BOB problem is 
significantly larger than the complexity of the strictly stochastic case.
Uniformly scaling the gaps down to satisfy $\gap_{(\nArms)}\leq1/8$
leaves all complexity ratios unchanged.
The square-root separation also persists, up to universal constants,
for the stochastic instances constructed in the proof of
Theorem~\ref{th:lowBOB} (Appendix~\ref{app:prooflowBOB}).
For sufficiently large $\nArms$, the optimal purely stochastic rate,
if guaranteed uniformly over stochastic instances and horizons, would
satisfy the premise of Theorem~\ref{th:lowBOB} at all sufficiently large
horizons. Hence it cannot be combined with vanishing worst-case error
over the adversarial class. Thus, in general,
the best performance achievable in the two worlds cannot be obtained
simultaneously.
	\end{remark}
%\begin{description}
%	\item[General Regime:] Let the set $S$ of pairs $(k,\alpha)$, be $S=\{k\in\setArmsmo,\alpha\geq 0 : \gap^{\alpha}_{(1)}=\frac{\gap^{\alpha}_{(k)}}{k}\}$. This set is not empty as for each $k$ it contains the pair $\left(k,\alpha=\frac{\log k}{\log\left(\frac{\gap{(k)}}{\gap_{(1)}}\right)}\right)$
%	Let us denote $\kappa=\max_{k,\alpha\in S} k^{\frac{1}{\alpha}}$
	%
	\paragraph{\raisebox{.04cm}{\textcolor{bull}{$\blacktriangleright$}}~Flat regime}  We assume all the gaps are equal, 
$k\in\setArms,~ \gap_k=\gap_{1}$. Then, 
$\complexitySR=\complexityBoth=\complexityUnif$. Having $\complexitySR=\complexityBoth$  shows that our  
stochastic BOB lower bound for robust learners (Theorem~\ref{th:lowBOB}) using $\complexityBoth$ is of the same order as the one in the 
strict stochastic setting~\citep{Audibert10BA} using $\complexitySR$.   
In this stochastic regime, \RULE{} is optimal while being robust to an adversary.
%
\paragraph{\raisebox{.04cm}{\textcolor{bull}{$\blacktriangleright$}}~Super-linear gaps} Let $2\in\argmin_{k\in\setArms} (\gap_{(k)}/k).$ 
This holds if 
$\gap_{(k)}=k\gap_{(1)}$ for every rank $k$ with $3\leq k\leq\nArms$,
and $\gap_{(1)}\leq 1/(8\nArms)$. Then, 
$\complexitySR=\complexityBoth=(2/\nArms) \complexityUnif$. 
Again, our BOB lower bound is of the same order as in 
the strict stochastic setting. This seems to indicate that in this case, BOB is achievable. However, it is not achieved by the uniform 
learner that is clearly suboptimal. This motivates the design of a learner that adapts its sampling strategy. Intuitively,  the learner can  
identify bad arms quickly
and  start focusing early on the best arms without incurring 
high variance on its estimates for them.
%
%\todov{think more about the right 2s!}
\paragraph{\raisebox{.04cm}{\textcolor{bull}{$\blacktriangleright$}}~Square-root gaps}   We assume $2\!\in\!\argmin_{k\in\setArms} (\gap^2_{(k)}/k).$ 
%This means that for all $k\in\setArms$, $\gap_{(1)}\leq \frac{\gap{(k)}}{\sqrt{k}}$. 
\newline
Choose a rank $j\in\argmin_{k\in\setArms} (\gap_{(k)}/k)$. 
For some constant $c$, let $\gap_{(1)}=c\gap_{(j)}/j$. We have 
$c\leq \sqrt{2j}$ because  $\gap^2_{(1)}/2 \leq \gap^2_{(j)}/j$ 
as  $2\in\argmin_{k\in\setArms} (\gap^2_{(k)}/k)$.
Therefore,  \[\complexitySR=\frac{2}{\gap^2_{(2)}}=\frac{2}{\gap^2_{(1)}}= \frac{2}{\gap_{(1)}} \frac{j}{c\gap_{(j)}}=\frac{2\complexityBoth}{c}=\frac{2\complexityUnif}{\nArms}\cdot\]	
We can get $c=\sqrt{2j}=\sqrt{2\nArms}$.
This happens if 
$\sqrt{\nArms/2}\gap_{(1)}=\gap_{(\nArms)}$ and
$\sqrt{k/2}\gap_{(1)}\geq\gap_{(k)}$ for ranks $k$ with $3\leq k<\nArms$. 
Then,  we get  $\sqrt{\nArms/2}\complexitySR=\complexityBoth$.
$\complexityBoth$ is $\sqrt{\nArms/2}$ larger than the complexity of
the strictly stochastic setting.
Intuitively,  the learner needs to spend some time to 
identify the `square-root gaps' suboptimal arms before starting 
to focus on the best arms. This makes it suffer an additional amount
of variance on its estimates for the best arms.
%\vspace{.1cm}
%
%\textbullet~ If all the gaps are such that for all $i$ , $\gap_i/i$ are
% equal (or more generally if, which means
% $\gap_1$ is quite small comparatively to the other gaps, the gaps in the example are increasing rapidly : linearly), we have the lowerbound 
%  $\complexitySR=\complexityBoth=\frac{1}{\nArms}\complexityUnif	$
%   the normal lowerbound again
%
%\textbullet~ If all the gaps are such that for all $i$ , $\gap^2_i/i$ are
% equal (or, more generally, If $1=\argmin_k \gap^2_k/k$, which 
% means arm 1 is quite small comparatively but as small as in
%  the previous case!, the gaps are increasing in a square root fashion ), in this case we have  $K=\argmin_k \gap_k/k$ we have the lowerbound 
%  $\complexitySR=\frac{1}{\sqrt{\nArms}}\complexityBoth=\frac{1}{\nArms}\complexityUnif	$ 
%  which is a new lowerbound larger than in the traditional case so 
%  thats cool. Indeed note that this is a case where the definition of the 
%  classical complexity is $\exp(-N\gap^2_1)$
%  
%  \todovout{\textbullet~ Think about other examples, what can we say in general? We can add something bout general geometric setting where all $\frac{i}{\gap^\alpha_i}$ are the same.}

%
%\begin{remark}
%The assumption $\nArms\exp\left(-\frac{1}{32}\gap^2_{(1)}\timeHorizon\right)\leq 1/32$ is mild. 
%Essentially, at most, it  asks for $\timeHorizon$ to be large 
%enough so that the stochastic problem is solvable within the budget $\timeHorizon$.
%% $\timeHorizon$ is required to be large enough here.
%\end{remark}%
%

\vspace{.5cm}
\noindent\textbf{Sketch of the proof} 
\textit{(full proof in 
	Appendix~\ref{app:prooflowBOB})}:
The idea is again to construct two environments that are difficult for the
learner to distinguish but have different best arms. This time, one environment
is stochastic and the other is adversarial. If the learner identifies the best
arm with very high probability in the stochastic environment, it cannot
reliably recognize that a different recommendation is needed in the
adversarial environment. Consequently, sufficiently good stochastic
performance forces a substantial probability of error in the adversarial
environment.

Our proof of the lower bound uses some arguments of purely stochastic best-arm identification 
lower bounds of~\cite{Audibert10BA} and~\cite{Carpentier16TB}.
We were also inspired by the lower bound of~\cite{Auer16AA} 
for the BOB question for the cumulative regret. However, 
our specific construction is new.
%\todov{thats a long sketch! should we shortening it or remove it? or no}

Fix the reference gaps, a learner, and a horizon $\timeHorizon$
satisfying Theorem~\ref{th:lowBOB}.
We construct a stochastic and an adversarial problem.
Between the two problems, only one arm differs.
Let us define $\nArms$ base Bernoulli distributions with means
$\mean_1^{\basePro}=1/2$ and
$\mean_j^{\basePro}=1/2-\gap_j$ for $j\in\setArmsmo$.
Choose any maximizing index
$i\in\arg\max_{j\in\setArms}j/\gap_j$; since $\gap_1=\gap_2$, $i\geq2$.
The initial phase is the prefix of $\timeHorizon_i$ rounds, where
\[
\dividefac_i=\min\left\{1,\frac{2\gap_1}{\gap_i}\right\},
\qquad
\timeHorizon_i=\left\lceil\timeHorizon\dividefac_i\right\rceil
\leq\timeHorizon.
\]
If $\timeHorizon_i=\timeHorizon$, as when $i=2$, the second phase is empty.

By Dirichlet's box principle, for $i\geq3$ at least one arm, denoted $\refarm$, among
$\lceil i/2\rceil,\ldots,i$ has expected pull count at most
$2\timeHorizon_i/i$. The maximizing property of $i$ makes its gap
comparable to $\gap_i$. For $i=2$, take $\refarm=2$. In both cases,
\[
\Exp_{\basePro}[\pullsNumber_{\refarm}(\timeHorizon_i)]
\leq\frac{2\timeHorizon_i}{i},\qquad
\frac{\gap_i}{2}\leq\gap_{\refarm}\leq\gap_i,\qquad\refarm\neq1.
\]
This lightly sampled arm is chosen from expectations under \basePro{},
before sampling the constructed tables.

The stochastic problem \stoPro{} is made by only modifying the original
Bernoulli distribution of arm $\refarm$, setting its mean to
$1/2+\gap_1/4$. Its distributions are stationary, and $\refarm$ is
uniquely best. The adversarial problem \advPro{} samples $\gainVector$ randomly
before interaction, independently of the learner's randomness.
It follows \basePro{} through the first $\timeHorizon_i$ rounds and
\stoPro{} afterward; all other arms keep their distributions throughout.
To make the resulting lower bound as strong as possible, we seek a
short initial phase to limit the distinguishing information.
However, it must be long enough to keep arm $1$ uniquely best on a
high-probability good-table event and preserve the prescribed gaps
after reward fluctuations. The chosen fraction provides this margin
while remaining of order $\gap_1/\gap_i$, up to the cap at one.
The appendix verifies that \stoPro{} and every good \advPro{} table
belong to $\Gaps_4$.

Let $P_A,P_S$ be the distributions of the learner's observed actions,
rewards, and final recommendation under \advPro{} and \stoPro{}.
To compare these observations, we use the change-of-measure
method of \citet*{Kaufmann16CBI}, in the form proved in
Lemma~\ref{l:transcriptChangeMeas}.
The lemma expresses the relative entropy of the two observation
distributions as the expected number of informative pulls times
the information supplied by each pull:
\[
D(P_A\Vert P_S)
=\Exp_{\basePro}[\pullsNumber_{\refarm}(\timeHorizon_i)]
\operatorname{kl}\left(\frac12-\gap_{\refarm},\frac12+\frac{\gap_1}{4}\right).
\]
Here the \advPro{} and \basePro{} interaction laws agree through the
prefix. The appendix gives, for universal constants $C,C'$,
\[
D(P_A\Vert P_S)
\leq\frac{C\timeHorizon_i\gap_i^2}{i}
\leq\frac{C'\timeHorizon\gap_1\gap_i}{i}
=\frac{C'\timeHorizon}{\complexityBoth^{\basePro}}.
\]
We want the lower bound on the learner's error to be as large
as possible. When the learner meets the prescribed guarantee
on \stoPro{}, a smaller upper bound on $D(P_A\Vert P_S)$ forces
a larger lower bound on its error in \advPro{}: its observations
provide less evidence for distinguishing the two problems.
This explains our choice of $i$: maximizing $i/\gap_i$ minimizes
the expression $\timeHorizon\gap_1\gap_i/i$ in the bound above.
For this maximizing index,
$\complexityBoth^{\basePro}=i/(\gap_1\gap_i)$, which gives the
dependence on $\timeHorizon/\complexityBoth^{\basePro}$.
The $\gap_j$ are the reference gaps; after reordering the constructed
\stoPro{} gaps, the appendix shows
$\complexityBoth^{\stoPro}\leq16\complexityBoth^{\basePro}$.
Thus the information is also bounded by a universal constant times
$\timeHorizon/\complexityBoth^{\stoPro}$.

Recommending arm $1$ is an error in \stoPro{}.
If its error satisfies Theorem~\ref{th:lowBOB}'s premise, this
recommendation is very unlikely under $P_S$. The lemma's event
consequence prevents it from being sufficiently likely under $P_A$
for reliable identification of arm $1$, forcing a positive constant
adversarial error. Concentration controls the bad tables, and the
appendix extracts a deterministic $\gainVector\in\Gaps_4$ with error
at least $1/16$.
The constructed instances may depend on $\timeHorizon$; this suffices
against guarantees required uniformly over instances and budgets.
%